% UC-071: Background and related work. Cite only entries verified in docs/literature.md.
% Sources marked "verified (metadata)" are cited only for what docs/literature.md records
% next to them (e.g. "as attributed by ...").
% TODO(verify): add verified entries for Zadeh's computing with words and Mendel's
% Perceptual Computer to docs/literature.md, then cite them in Section 2.1.
% Pincay 2022 (issue #53 open): claim only what pincay2021 §5 states.
\section{Background and Related Work}
\label{sec:related}

The assistant draws on four strands of work: computing with words and its use in citizen
interfaces, fuzzy linguistic summaries and their quality measures, explanation through
linguistic summaries, and the verification of quantified statements produced by language
models. A fifth, smaller strand concerns how passengers perceive punctuality.

\subsection{Computing with Words in Citizen-Facing Systems}
\label{sec:rw-cww}
Computing with words (CWW) treats words as the objects of computation: their meaning is a
fuzzy set, and a system reasons with these meanings instead of with crisp categories. A
common way to structure such a system is the architecture of the Perceptual Computer, in
which an \emph{encoder} turns inputs into fuzzy sets, a \emph{CWW engine} computes with
them, and a \emph{decoder} turns the result back into words a person can act on. Our
assistant follows this structure (Section~\ref{sec:concept}).

At Human-IST, CWW has been studied as a way to make interfaces respond to people's
perceptions rather than to rigid categories: \citet{colombo2024} proposes \emph{phenotropic
interaction}, interfaces improved with computing with words and perceptions. The limits of
crisp categories in public data are the subject of \citet{spycher2026}, who examine
taxonomies for open government data. Our work shares this perspective: the official
punctuality rule is a crisp category over open government data, and a passenger's notion of
``late enough to miss my train'' is a perception that depends on their situation.

\subsection{Fuzzy Linguistic Summaries and Their Quality}
\label{sec:rw-ls}
A linguistic summary is a quantified sentence about a data set, such as ``most trains are on
time''. The idea goes back to \citet{yager1982}. \citet{kacprzyk2009} describe such
summaries through \emph{protoforms}, templates such as ``$Q$ of the $x$ are $S$'' and ``$Q$ of
the $R$ $x$ are $S$'', in which $Q$ is a fuzzy quantifier (``most''), $S$ a summarizer
(``on time'') and $R$ a restriction (``in the morning peak''). The truth, or validity, of a
summary is the membership of the summarizer's proportion in the quantifier; we use it in
the form given by \citet[Eq.~(2)]{pincay2021}, who attribute it to \citet{kacprzyk2009}, and
by \citet[Eqs.~(5)--(6)]{hudec2022}, who attribute it to \citet{yager1982}.

Validity alone does not say how much data a sentence rests on. \citet{hudec2018} use short
quantified sentences to disseminate official statistics and combine validity with
\emph{coverage}: a restricted summary may be fully valid and yet describe only a handful of
records. Their quality measure $Q_c$ is a non-idempotent t-norm of validity and coverage
(for example the product, which we use) when the coverage reaches a threshold, and zero
otherwise \citep[Eq.~(9)]{hudec2018}; they attribute
the merge of validity and coverage to \citet{hudec2017}. The degree of sufficient coverage is
a smooth function of the share of records that support both the restriction and the
summarizer \citep[Eqs.~(A.3)--(A.4)]{hudec2018}, a function \citet{hudec2022} attribute to
\citet{wu2010}. \citet{hudec2022} further propose a \emph{regional} quality measure for
summaries shown on maps: the same sentence is evaluated per region, and its truth is combined
with the region's relative size through an ordinal sum of a conjunctive and a disjunctive
operator, so that small regions are neither ignored nor allowed to dominate. These are the
measures we use for SRQ1.

\subsection{Linguistic Summaries of Traffic Data}
\label{sec:rw-traffic}
\citet{pincay2021} mine linguistic summaries from road traffic data with the quantifiers
\emph{few}, \emph{about half} and \emph{most}, a validity threshold for reporting a summary,
and the Simpful library for fuzzy logic \citep{spolaor2020}. In their conclusion they name
two extensions as future work: combining validity with existing quality measures in the
sense of Hudec et al., and checking validity by region, to see whether the data distribution
skews the summaries \citep[§5]{pincay2021}. The same group has continued with geospatial
traffic data \citep{pincay2022} and with Swiss traffic safety analysis combining machine
learning and a local language model \citep{pincay2026}. Our work transfers the approach to
public transport and implements both extensions named in \citet{pincay2021}; we use their
configuration as a methodological baseline, not as a reproduction of their numbers, since
their data is not public.

\subsection{Linguistic Summaries as Explanations}
\label{sec:rw-xai}
Linguistic summaries have been proposed as a form of explainable artificial intelligence.
\citet{plenary2022} explain the behaviour of black-box models in natural language through
fuzzy linguistic summaries, and \citet{wrede2023} work towards making linguistic summaries
themselves explainable. In PLENARY the summaries describe a separate black-box model; in our
assistant the summaries are the model: the decision word is decoded from
the same fuzzy proportion that the main explanation sentence states, so the explanation
cannot diverge from how the decision was made.

\subsection{Language Models and Quantified Statements}
\label{sec:rw-llm}
\citet{eida2026} study how language models generate quantified statements from tables and
verify them by executing checks against the source table. Their truth conditions are crisp
and set-based; ``most'' is true if more than half of the relevant rows satisfy the property
\citep{eida2026}. Their tables are synthetic and small, and their verification
functions are themselves written by a language model. We adopt their set-based truth
conditions as our crisp checker, implemented as deterministic code, and add a fuzzy checker
that assigns graded truth to vague quantifiers and summarizers, on real data.
\citet{gaitan2025} combine language models with fuzzy linguistic summaries of raw sensor
data in diabetes care, with a methodology for fine-tuning and evaluating the models; the
truth of the quantifiers that the language model writes is not recomputed from the data. The
language model in our audit is Apertus \citep{apertus2025}, an open model from the Swiss AI
Initiative that we run locally.

\subsection{Passengers' Perception of Punctuality}
\label{sec:rw-perception}
Official punctuality measures and passengers' experience differ. A study of passengers in
Great Britain found that passengers judge punctuality more strictly than the industry
measure and expect trains to arrive within about one minute of the timetable
\citep{transportfocus2015}. We use this finding as motivation only. Swiss perception data
does not exist for our purpose, and we do not collect new data from people. Instead, our
perception-aware view makes ``late'' depend on a quantity the passenger knows precisely:
their own time to change trains.

\subsection{Positioning}
\label{sec:rw-position}
Taken together, linguistic summaries with quality measures exist for official statistics
and maps \citep{hudec2018,hudec2022}, linguistic summaries of traffic exist without them
\citep{pincay2021}, linguistic summaries explain other models \citep{plenary2022} and are
themselves made explainable \citep{wrede2023},
and quantified statements of language models are verified crisply \citep{eida2026}. This
paper connects these strands in one citizen-facing artefact: summaries with Hudec's quality
measures on public transport data are the decision core of an assistant, the decision depends
on the passenger's own margin, and the same truth computation serves as a fuzzy checker for
language-model answers, compared against a crisp one.
